\section{Research context and the results}
\label{sec:context}

The useful common ground between the two repositories is finite Fourier
analysis with explicit verification. The OpenAI manuscript
\emph{Unbounded Violations of the Square-Root Degree Bound}
\cite{OpenAI2026} develops low-degree observations whose conditional
coordinate sums retain unexpectedly large variance. ProveIt's
\emph{Noisy parity cubes} report \cite{ProveItNoise} studies spectral
samples, parity noise, and query costs, while its Thue--Morse Lean files
\cite{ProveItWalsh} establish finite character identities. This article
uses the former observation mechanism, develops a new structural
threshold and rigidity theory, and follows the latter repository's
practice of separating written proofs from finite computational checks.

\subsection{What was reviewed}

The OpenAI repository catalog and overview were examined to choose a
focused problem. The proof-level review for this article covers all seven
section files of result family 192: the observation calculus, reporting
construction, alternative reports, mean-cost extension, weighted
Gaussian mixture, and small-width calculation. The fixed source revision
is
\[
\texttt{adc7f1241b42e322a6451854ab7e4b4c146bf78a}.
\]
The ProveIt review covers the actual source of \emph{Noisy parity cubes},
its verification script, and the two Lean files
\texttt{ThueMorseBooleanCube.lean} and \texttt{ThueMorseWalsh.lean}, at
revision
\[
\texttt{dfc1a5a6817462e0465e445e1466c41ee42ea839}.
\]
The repository README was also inspected at the later snapshot
\texttt{1aa2f12ec85b}.
These are scoped mathematical reviews, not audits of every result in
either repository. Exact paths and source identities are recorded in the
accompanying provenance file.

Theorem 3.1 of \emph{Noisy parity cubes} represents a Boolean rule's
two-input disagreement under independent XOR noise as a mixture of the
corresponding parity disagreement probabilities, with squared Walsh
coefficients as weights. Its Lemma 3.2 compares total influence with
expected exact query cost. Those conclusions concern noise response and
queries; they do not imply a variance bound for arbitrary low-degree
cells. The inspected Lean files prove product identities and vanishing
nontrivial Walsh-character sums in a binary representation. They do not formalize the
conditional-expectation results proved here.

\subsection{The distinction between three Fourier quantities}

For \(f:\{-1,1\}^N\to\{-1,1\}\), write
\[
 L(f)=\sum_{i=1}^N\widehat f(\{i\}),\qquad
 L_{\mathrm{abs}}(f)=\sum_{i=1}^N|\widehat f(\{i\})|.
\]
These differ from the squared first-level Fourier weight and from total
influence. Coordinate sign changes preserve degree and turn the signed
sum into the absolute sum. Thus their universal extremal inequalities
are equivalent, although the values for a particular function need not
be equal.

The Gopalan--Servedio problem asks whether \(L(f)\le\sqrt{\deg f}\);
O'Donnell's problem list records it and the stronger majority benchmark
\(L(f)\le B_{\deg f}\), where \(B_d=\E|H_d|\)
\cite{ODonnell2012}. The same square-root question appears as
Conjecture 3.17 in \cite{FHKL2016}. The source manuscript
\cite{OpenAI2026} proves violations by every constant factor. The
present article does not claim that prior conclusion as a new theorem.

For coordinate-subcube partitions, the retained variance equals the
expected number of fixed coordinates. This is the conditional-sum
calculation behind the decision-tree inequality of O'Donnell and
Servedio \cite[Theorem 3]{ODonnellServedio2007}. Blais, Tan, and Wan
\cite[Lemma 3.1]{BlaisTanWan2015} use the same type of conditional sum
for parity decision trees. General low-degree indicators allow
cancellation unavailable to coordinate subcubes. The codimension
question asks how far that cancellation can first go.

\subsection{The structural threshold}

Let \(D(F)\) denote maximum cell degree, as defined precisely in
Section~\ref{sec:observations}. The principal result is the following.

\begin{theorem}[First possible codimension for a variance gain]
\label{thm:headline}
Let \(X_1,\ldots,X_N\) be independent uniform signs and
\(H_N=\sum_iX_i\).
\begin{enumerate}
\item If \(C\ne\varnothing\) and
\(\deg\1_C\le N-k\), where \(k\in\{1,2,3\}\), then
\[
 \Var(H_N\mid C)\ge k.
\]
Equality holds if and only if \(\Law(H_N\mid C)\) is an integer translate
of the law of a sum of \(k\) independent uniform signs.
\item Consequently, if \(D(F)\le N-k\), then \(v(F)\le N-k\).
Equality holds if and only if every attained cell has the equality
distribution in the first assertion.
\item There is an explicit twenty-bit observation \(F_*\) with
\[
 D(F_*)=16,\qquad
 v(F_*)=16+\frac9{34816}.
\]
Its exceptional cell \(A\) has exact degree 16 and
\[
 \Var(H_{20}\mid A)=4-\frac9{289}.
\]
\end{enumerate}
Thus four is the smallest Fourier codimension \(N-D(F)\) at which
retained variance can exceed cell degree.
\end{theorem}

The first two assertions are proved using interlacing even and odd
integer lattices. The last is a completely specified finite example,
not a limiting construction. Its numerical excess is small, but it
locates an exact boundary that the qualitative amplifier does not
identify.

The equality statement is distributional. It does not classify a cell
as a coordinate subcube. Indeed, the proof only requires certain
alternating moments of the Hamming-weight distribution to vanish.
The stability theorem makes this distinction quantitative. If
\(\Delta=(N-k)-v(F)\), then for suitable integer shifts \(c(F)\) and an
independent \(k\)-sign sum \(S_k\),
\begin{equation}\label{eq:headline-stability}
 \TV\bigl(\Law(F,H_N),\Law(F,c(F)+S_k)\bigr)
 \le C_k\Delta,\qquad
 (C_1,C_2,C_3)=\left(\frac14,\frac14,\frac38\right).
\end{equation}
We also prove sharper local bounds and construct extremizers in the
abstract moment class.

\subsection{Other contributions and their scope}

The finite reporter is reduced to an exact rational optimization
formula, and its full score law is computed as a four-atom perturbation
of \(H_{16}\). Tensorization proves that no universal additive constant
can repair \(v(F)\le D(F)\). Exhaustive rational enumeration proves
twenty-bit minimality within the family of equal binomial blocks,
sum-level leaf tests, and a central-sum selector. This is not a
classification of every observation on fewer than twenty bits.

For the source Gaussian amplifier, we prove explicit restricted-moment
estimates, a uniform fourth-moment invariant, and a single copy count
valid at every iteration. The resulting construction has polynomial
dimension growth in its prescribed degree budget. A fully rational
choice supplies the numerical exponent
\(\varepsilon_0=2^{-8\,000\,000\,000}\).
Its size makes the construction an effective existence theorem rather
than a practical way to tabulate a counterexample.

There is a necessary priority distinction here. Once an arbitrary-factor
square-root violation is known, balanced block composition already
implies a positive power separation. Fact 2.7 of
\cite{BlaisTanWan2015} gives the relevant multiplication rule.
Section~\ref{subsec:composition} proves the short deduction, including
the balancing step. The additional content of our analytic construction
is the fixed quantitative schedule and explicit constants, not the bare
existence of a positive exponent.

All main mathematical claims are proved below. Their priority relative
to the entire literature remains a matter for further scholarly review;
the comparison here is with the repository sources and primary papers
actually inspected.
